\subsection{Наблюдения}

\begin{itemize}
\item \textbf{Москва --- половина рынка и самые высокие зарплаты.} 26\,726 из 54\,607 вакансий варианта (49\%). Медиана 150~тыс. на руки --- на 36\% выше Петербурга (110~тыс.) и на 88\% выше Екатеринбурга (80~тыс.). Даже нижний квартиль Москвы (100~тыс.) выше медианы Екатеринбурга.
\item \textbf{Разрыв объясняется и составом вакансий, и ценой навыка.} В Москве больше вакансий для опытных: 46\% с опытом от 3 лет против 38\% в Петербурге и 26\% в Екатеринбурге; без опыта --- 6\%, 9\% и 11\%. Но и за \emph{одинаковый} навык Москва платит больше: SQL --- 174 / 125 / 104~тыс., системный анализ --- 200 / 150 / 120~тыс., UML --- 218 / 150 / 138~тыс. Надбавка Москвы за тот же навык --- от трети до двух третей.
\item \textbf{Удалёнка не выравнивает рынок.} Доля удалённых вакансий почти одинакова (24--25\% во всех трёх регионах), но зарплаты по-прежнему привязаны к городу публикации.
\item \textbf{Разный спрос на навыки.} В Москве и Петербурге первый навык --- SQL (37,6\% и 30,5\% вакансий с навыками), в Екатеринбурге --- «аналитическое мышление» (29,4\%), SQL только второй. Тройка BPMN + SQL + UML --- самая частая во всех трёх регионах, то есть ядро системного аналитика общее; различаются окраины.
\item \textbf{Характерные навыки показывают профиль рынка.} Москва --- большие данные и хранилища: Hadoop (в 3,3 раза чаще, чем в двух других регионах), Greenplum (3,1×), DWH (2,1×) --- банки и крупные IT-компании. Петербург --- внедрение 1С ERP: 1С:УП, 1С:КА, 1С:ERP, 1С:УПП (2,2--2,5×) --- производственные и торговые компании. Екатеринбург --- общие формулировки («аналитический склад ума», «работа в команде», «информационные технологии») и программирование 1С: вакансии менее специализированы, что согласуется с самой низкой медианой.
\item \textbf{Чем не различаются.} Доля вакансий с навыками (60--65\%) и длина списка навыков (6,7--6,8) почти одинаковы: регионы различаются не количеством требований, а их содержанием и ценой.
\end{itemize}
